\section{Introduction}

Video generation and understanding represents one of the most challenging frontiers in artificial intelligence, requiring models to comprehend temporal dynamics, physical causality, and visual coherence across multiple modalities. This literature review surveys the leading academic research on "Video Models" from elite university laboratories, examining the architectural innovations, training methodologies, and evaluation benchmarks that have propelled this field forward.

The research landscape has been transformed by the convergence of transformer architectures with diffusion models, culminating in systems capable of generating minute-long videos with unprecedented fidelity. This review identifies five primary research thrusts: (1) latent video tokenization and compression, (2) diffusion transformer architectures for video generation, (3) world models and physical commonsense learning, (4) efficient and scalable video generation systems, and (5) video understanding through self-supervised reasoning.

We systematically analyze the contributions from leading university laboratories including UC Berkeley's Berkeley Artificial Intelligence Research (BAIR) lab, MIT's HAN Lab, Stanford University's Human-Centered Artificial Intelligence (HAI) institute, and Carnegie Mellon University's Robotics Institute. Each laboratory has made seminal contributions that have shaped the trajectory of the field.

\section{Top University Laboratories}

\subsection{University of California, Berkeley - BAIR Lab}

The Berkeley Artificial Intelligence Research (BAIR) laboratory has been a pioneer in video generation, particularly through the work of Pieter Abbeel's group. Wilson Yan's doctoral dissertation "Learning About The World Through Video Generation" represents a comprehensive investigation into the fundamental challenges of scaling generative models to natural videos.

The BAIR lab's research program has systematically addressed computational bottlenecks in video generation through three major innovations: VideoGPT (2021), TECO (2023), and ElasticTok (2024). This sequence of work demonstrates a principled approach to reducing computational requirements while maintaining generation quality. The VideoGPT architecture introduced a VQ-VAE based framework that compresses video into discrete latent representations, enabling transformer-based autoregressive generation at scale.

\textbf{Influence Index: 92/100}

The BAIR lab's work has received significant recognition, with VideoGPT accumulating over 800 citations. The research has directly influenced subsequent developments in latent video diffusion models and has established architectural patterns adopted by both academic and industrial laboratories worldwide.

\subsection{MIT HAN Lab}

The MIT HAN Lab, led by Associate Professor Song Han, has focused on efficient video generation through innovative architectural designs. The laboratory's research emphasizes practical deployment, developing models that can run on consumer hardware while achieving competitive quality.

The SANA-Video system (2025) introduces block linear attention as a replacement for standard quadratic self-attention, enabling efficient generation of minute-long videos at 720×1280 resolution. This work demonstrates that high-quality video generation need not require massive computational resources, achieving training costs仅为 1\% of comparable industrial systems.

The laboratory has also contributed to the broader diffusion transformer ecosystem through SANA (2024), which introduced linear attention mechanisms for high-resolution image synthesis, and HART (2025), which presents a hybrid autoregressive-diffusion approach for efficient visual generation.

\textbf{Influence Index: 88/100}

MIT HAN Lab's work has been recognized through multiple publications at top venues including ICLR (oral), CVPR, and ICML. The practical focus on efficiency has influenced both academic research directions and industrial deployment strategies.

\subsection{Stanford University - HAI and Computer Science}

Stanford's research contributions span multiple aspects of video models, from generation to understanding. The Video-STaR system (2025), developed in collaboration with Google Research, introduces self-training approaches that enable models to reason about video content without requiring extensive human annotations.

The Unified Video Action Model (UVA), developed at Stanford, addresses the integration of video generation with action prediction for robotics applications. This work demonstrates that video models can serve dual purposes—as generative systems and as world models for planning and control.

Stanford researchers have also contributed critical analysis of video generation capabilities, including the influential "How Far is Video Generation from World Model" study that evaluates physical commonsense understanding in video generation systems.

\textbf{Influence Index: 85/100}

Stanford's location within Silicon Valley has enabled close collaboration with leading technology companies, facilitating the transition of academic innovations to practical applications. The university's interdisciplinary approach, exemplified by the HAI institute, ensures that video model research considers both technical and societal implications.

\subsection{Carnegie Mellon University - Robotics Institute}

Carnegie Mellon University's Robotics Institute has made foundational contributions to video tokenization and efficient generation. The MAGVIT system (2022) introduced the masked generative video transformer, achieving state-of-the-art results across multiple video generation benchmarks while providing 60× speedup over autoregressive alternatives.

More recently, the UniDisc system (2025) presents a unified multimodal discrete diffusion model capable of generating both images and text through a single framework. This work represents an important direction toward truly general-purpose generative models that can handle multiple modalities within a unified architecture.

CMU researchers have also contributed to physical AI systems that leverage video generation for robot manipulation and scene understanding.

\textbf{Influence Index: 82/100}

CMU's long-standing expertise in robotics has naturally extended to video models as potential world models for embodied agents. The university's substantial infrastructure, including the NVIDIA partnership for GPU computing, enables research at scale.

\section{Key Papers and Citation Analysis}

\subsection{VideoGPT: Video Generation using VQ-VAE and Transformers}

\textit{Authors: Wilson Yan, Yunzhi Zhang, Pieter Abbeel, Aravind Srinivas (UC Berkeley)}

\textit{Citations: 807+}

This foundational paper introduced the VideoGPT architecture, demonstrating that autoregressive transformer models could be effectively applied to video generation through VQ-VAE latent compression. The work established a template for subsequent latent video generation systems and demonstrated competitive results against GAN-based approaches on standard benchmarks including BAIR Robot Pushing, UCF-101, and Tumblr GIFs.

The paper's influence extends through its architectural choices—axial attention for efficient spatio-temporal computation and discrete latent representation—that have become standard in the field.

\subsection{Temporally Consistent Transformers for Video Generation (TECO)}

\textit{Authors: Wilson Yan, Danijar Hafner, Stephen James, Pieter Abbeel (UC Berkeley, DeepMind)}

\textit{Conference: ICML 2023}

TECO addresses the challenge of generating long-duration videos with temporal consistency. The paper introduces a latent dynamics model that separates spatial and temporal computation, enabling efficient scaling to sequences containing hundreds of frames. The work established new benchmarks on challenging datasets including Minecraft and Habitat environments, demonstrating generation of temporally coherent videos over extended time horizons.

\subsection{MAGVIT: Masked Generative Video Transformer}

\textit{Authors: Lijun Yu, Jose Lezama, Alexander G. Schwing (CMU, Google Research)}

MAGVIT introduced the masked token modeling approach to video generation, demonstrating that a single model could handle multiple generation tasks including class-conditional generation, video inpainting, and future prediction. The system achieves 60× speedup over prior autoregressive approaches while maintaining competitive quality.

\subsection{SANA-Video: Efficient Video Generation with Block Linear Diffusion Transformer}

\textit{Authors: Junsong Chen et al. (NVIDIA, MIT, Tsinghua)}

\textit{Conference: ICLR 2025 (submitted)}

This paper presents a complete system for efficient video generation that can run on consumer GPUs including the RTX 5090. The key innovation is block linear attention that maintains constant memory requirements regardless of video length, enabling minute-long generation with standard hardware.

\section{Citation Tree Analysis}

The citation tree for video generation models traces a clear lineage from early autoregressive models through latent diffusion approaches to the current diffusion transformer paradigm:

\begin{enumerate}
    \item \textbf{Root}: GAN-based video generation (2016-2020)
    \item \textbf{Generation 1}: VideoGPT (Yan et al., 2021) - VQ-VAE + Transformer
    \item \textbf{Generation 2}: MAGVIT (Yu et al., 2022) - Masked Transformer
    \item \textbf{Generation 3}: Stable Video Diffusion (Blattmann et al., 2023) - Latent Diffusion
    \item \textbf{Generation 4}: DiT-based models (2024) - Sora, CogVideoX
    \item \textbf{Generation 5}: Efficient DiT (2025) - SANA-Video, LongLive
\end{enumerate}

Each generation has built upon the innovations of previous work while introducing architectural innovations that address the fundamental challenges of video generation: temporal consistency, computational efficiency, and long-range coherence.

\section{Conclusions}

The field of video models has undergone rapid transformation, with elite university laboratories making foundational contributions that have shaped both architectural choices and application directions. Key themes that emerge from this review include:

\begin{enumerate}
    \item \textbf{Latent Space Efficiency}: The progression from pixel-space to latent-space generation has been crucial, enabling transformers to handle the combinatorial complexity of video data.
    \item \textbf{Architecture Unification}: The diffusion transformer paradigm has emerged as the dominant approach, replacing specialized architectures with general-purpose transformer designs.
    \item \textbf{Efficiency as Priority}: University research has emphasized practical efficiency, with significant contributions to linear attention, quantization, and training strategies that reduce computational requirements.
    \item \textbf{World Model Integration}: Video generation is increasingly viewed not merely as a generation task but as a path toward world models that can simulate physical environments.
\end{enumerate}

Future research directions emerging from this literature include improved physical commonsense reasoning in generated videos, more efficient long-context attention mechanisms, and tighter integration between video generation and robotic planning systems.

\section{Acknowledgments}

This literature review synthesizes contributions from multiple university laboratories. The author acknowledges the open sharing of research through arXiv and conference publications that enables comprehensive review of the field.